\subsection{Сопоставление с экспертным направлением}
\label{sec:ch5_analytical_expert_comparison}
\conceptsource{Статья 5, таблица 10 и пояснения
валидационной сетки. Сохранены разные n направления
при общей исходной сетке.}

\iffalse
(Примечание автора: Чё? Так мало? Всего лишь? Эмм… Так может нам перезапустить на большем числе размеченных данных? Ведь уже есть.. Хотя понял, валидационные. Но погодь, можно же все, зачем тут только валидационные? Мы же не нейронку проверяем которую учили. Ну можно и разбить данные на «всё» и «валидационное».)
\fi
\authorreviewnote{AA178}

В исходном экспертном сопоставлении использованы
54 валидационные записи, из них 32 Open\_EE
и 22 French/RTE.
Общая сетка содержит 12\,942 точки применимости,
в 10\,181 из которых эксперт определил направление.
Внутри файла применяется шаг пять отсчётов
и ограничение 256 точками.
Это уменьшает преобладание длинных записей,
но не превращает соседние точки
в независимые аварийные события.

Начальное накопление контекста и переходная
маска в этой сетке исключены.
Направление оценивается при совместной применимости
эксперта и проверяемого органа.
Поэтому фактическое число сравниваемых точек
различается, как показано в
таблице~\ref{tab:ch5_analytical_expert}.
Неприменимые ответы оцениваются отдельно,
а не засчитываются правильным направлением.

\begin{table}[htbp]
\centering\small
\caption{Аналитические РНМ на исходной экспертной
валидационной сетке}
\label{tab:ch5_analytical_expert}
\begin{tabularx}{\textwidth}{@{}Xrrr@{}}
\toprule
Орган & Точек направления & Macro-F1, \% & MCC \\
\midrule
Адаптивный & 10\,176 & 89,59 & 0,794 \\
Пофазный угловой & 10\,008 & 44,79 & 0,037 \\
Угловой прямой последовательности & 9998 & 58,17 & 0,231 \\
Пофазный мощностной & 10\,008 & 46,10 & 0,159 \\
Мощностной прямой последовательности & 9998 & 55,57 & 0,279 \\
Двухвыборочный & 10\,013 & 52,77 & 0,153 \\
Пятивыборочный & 10\,089 & 58,73 & 0,325 \\
БМРЗ, дополнительная ветвь & 9998 & 65,62 & 0,484 \\
Адаптация БАВР-072 & 10\,157 & 58,68 & 0,241 \\
\bottomrule
\end{tabularx}
\end{table}

\iffalse
(Примечание автора: И непонятно что такое MCC, если не описывалось где-то ,может забыл.)
(Примечание автора: Ок, закрываем задачу.)
\fi
\authorreviewnote{AA179}

MCC учитывает четыре исхода двоичного сравнения.
Единица означает полное совпадение,
ноль отсутствие выявленной связи ответов с эталоном.
Macro-F1 одинаково учитывает два направления.
Оба показателя описывают принятую сетку
и правила допуска к сравнению.
Они дополняются проверкой полного состояния
и ошибочного времени из
раздела~\ref{sec:ch5_engineering_verification}.

На этой сетке адаптивный орган имеет
наибольшее согласие направления с экспертом.
Но объединённый результат скрывает различия источников.
Для пятивыборочного органа Macro-F1
составляет 94,67\% во French/RTE
и 28,67\% в Open\_EE.
Для БМРЗ-варианта получены 89,84 и 46,89\%.
Такое расхождение важно для выбора
направления модернизации применительно
к промышленному БАВР.
Оно требует разбора режимов и целевой логики,
а не только выбора максимального числа
в объединённой таблице.

Эта проверка показывает, как массовое сравнение
переходит к предметной оценке конкретного органа.
Высокое согласие в одном источнике
не освобождает от проверки другого,
а изменение формулы проверяется
по тем же требованиям к выходу.
Дальнейшее сравнение с обучаемыми моделями
сохраняет эти единицы оценки и ограничения.


